\documentclass[11pt,letterpaper]{article}
\usepackage{../../Common/course_notes}
\usepackage{../../Common/course_macros}

\begin{document}
\setlength{\abovedisplayskip}{6pt}
\setlength{\belowdisplayskip}{6pt}
\setlength{\textfloatsep}{10pt}
\setlength{\intextsep}{8pt}
\color{blackmain}

\begin{center}
{\LARGE\textbf{Geometría y Física I: Mecánica Clásica y Formas de Cartan}}\\[0.25cm]
{\Large\textcolor{orangeacc}{\textbf{Semana 1: Notación de Einstein, Delta de Kronecker y Símbolo de Levi-Civita}}}\\[0.15cm]
{\large\textbf{El fin del álgebra vectorial engorrosa: demoliendo identidades en 3 líneas}}\\[0.3cm]
{\large\textcolor{graymuted}{Notas de Estudio --- Nivel Licenciatura}}\\[0.2cm]
\end{center}

\vspace{0.2cm}

% =====================================================================
\manualsection[1. La Crisis del Cálculo Vectorial]{1. La Crisis del Cálculo Vectorial Tradicional}

En los cursos tradicionales de cálculo de una y varias variables, el álgebra vectorial se enseña casi exclusivamente mediante los vectores unitarios cartesianos $\{\hat{i}, \hat{j}, \hat{k}\}$ y la regla mnemotécnica del determinante $3 \times 3$ para el producto cruz:
\[
\vec{A} \times \vec{B} = \det \begin{pmatrix} \hat{i} & \hat{j} & \hat{k} \\ A_x & A_y & A_z \\ B_x & B_y & B_z \end{pmatrix} = (A_y B_z - A_z B_y)\hat{i} - (A_x B_z - A_z B_x)\hat{j} + (A_x B_y - A_y B_x)\hat{k}.
\]
Aunque este esquema resulta familiar para calcular torques o momentos angulares simples, se convierte en una auténtica pesadilla cuando necesitamos demostrar identidades físicas fundamentales.

\begin{intuitionbox}{El drama de la identidad BAC--CAB}
Consideremos la célebre identidad del doble producto cruz:
\[
\vec{A} \times (\vec{B} \times \vec{C}) = \vec{B}(\vec{A}\cdot\vec{C}) - \vec{C}(\vec{A}\cdot\vec{B}).
\]
Si intentamos demostrar esta igualdad por fuerza bruta expandiendo las tres componentes cartesianas, debemos calcular primero el vector intermedio $\vec{D} = \vec{B}\times\vec{C}$, luego cruzarlo con $\vec{A}$ componente a componente, y finalmente sumar y restar términos artificiales como $A_x B_x C_x$. El cálculo suele tomar **dos páginas completas de álgebra propensa a errores de signos**.

¿Por qué algo tan geométrico como orientar planos en el espacio requiere un esfuerzo algebraico tan ciego? La razón es clara: el formalismo tradicional nos obliga a escribir tres veces la misma ecuación (una para $x$, una para $y$, una para $z$) sin explotar la simetría de permutación del espacio tridimensional.
\end{intuitionbox}

En esta primera semana construiremos las tres herramientas fundamentales que transformarán este cálculo agotador en una deducción transparente de **tres renglones**:
\begin{enumerate}
    \item El \term{convenio de suma de Einstein}, que elimina la notación redundante de sumatorias.
    \item La \term{delta de Kronecker} $\delta_{ij}$, que actúa como un filtro e interruptor algebraico.
    \item El \term{símbolo de Levi-Civita} $\epsilon_{ijk}$, que codifica la orientación y la alternancia geométrica del producto cruz.
\end{enumerate}

% =====================================================================
\manualsection[2. Notación de Einstein]{2. El Convenio de Suma de Einstein y los Índices Libres vs. Mudos}

El primer paso consiste en unificar los nombres de los ejes espaciales. En lugar de escribir $(x, y, z)$, denotamos las coordenadas como $(x_1, x_2, x_3)$ o genéricamente $x_i$, donde el índice latino $i \in \{1, 2, 3\}$. De igual forma, las componentes de un vector $\vec{A}$ se escriben $(A_1, A_2, A_3)$.

Bajo esta convención, el producto punto euclidiano tradicional entre dos vectores es:
\[
\vec{A}\cdot\vec{B} = A_1 B_1 + A_2 B_2 + A_3 B_3 = \sum_{i=1}^3 A_i B_i.
\]
En 1916, mientras formulaba la Relatividad General, Albert Einstein notó que escribir el símbolo $\sum$ era completamente redundante: cada vez que calculamos una contracción física, el índice de la sumatoria aparece exactamente **dos veces**.

\begin{theorembox}{Convenio de Suma de Einstein}
Si un índice literal aparece exactamente \textbf{dos veces} en un mismo término monomio, se sobreentiende una sumatoria sobre todo el rango de valores de dicho índice (típicamente de $1$ a $3$ en el espacio euclidiano):
\[
A_i B_i \equiv \sum_{i=1}^3 A_i B_i = A_1 B_1 + A_2 B_2 + A_3 B_3.
\]
\end{theorembox}

\manualsubsection{Índices Mudos vs. Índices Libres}
Para operar sin errores en cálculo tensorial, es indispensable distinguir dos clases de índices:
\begin{itemize}
    \item \term{Índice Mudo (o de suma):} Es aquel que aparece dos veces en un término. Es una variable interna muda que puede renombrarse arbitrariamente sin alterar el valor de la expresión:
    \[
    A_i B_i = A_k B_k = A_m B_m = \vec{A}\cdot\vec{B}.
    \]
    \item \term{Índice Libre:} Es aquel que aparece exactamente **una vez** en cada término de una ecuación. Representa la componente vectorial que sobrevive tras las sumas:
    \[
    v_i = A_i + B_i \quad \Longleftrightarrow \quad \begin{cases} v_1 = A_1 + B_1 \\ v_2 = A_2 + B_2 \\ v_3 = A_3 + B_3 \end{cases}
    \]
\end{itemize}

\begin{infobox}[pureblack]{Las Dos Reglas de Oro de los Índices}
\begin{enumerate}
    \item \textbf{Conservación de índices libres:} El número y nombre de los índices libres debe ser exactamente el mismo a ambos lados de una igualdad. La expresión $v_i = A_j$ carece de sentido matemático.
    \item \textbf{Prohibido repetir más de dos veces:} Un índice jamás puede aparecer tres o cuatro veces en un mismo monomio. Una expresión como $A_i B_i C_i$ es inválida. Si necesitas sumar sobre varios vectores, debes usar letras distintas: $(A_i B_i)(C_j D_j)$.
\end{enumerate}
\end{infobox}

% =====================================================================
\manualsection[3. Delta de Kronecker]{3. La Delta de Kronecker $\delta_{ij}$: El Filtro de Sustitución}

¿Cómo escribimos la matriz identidad o un operador de proyección en notación de índices? Introducimos la \term{delta de Kronecker}:

\begin{theorembox}{Definición de la Delta de Kronecker}
Para cualesquiera índices $i, j \in \{1, 2, 3\}$, se define:
\[
\delta_{ij} = \begin{cases} 1 & \text{si } i = j, \\ 0 & \text{si } i \neq j. \end{cases}
\]
Matricialmente, $\delta_{ij}$ son las entradas de la matriz identidad $3 \times 3$:
\[
(\delta_{ij}) = \begin{pmatrix} 1 & 0 & 0 \\ 0 & 1 & 0 \\ 0 & 0 & 1 \end{pmatrix}.
\]
\end{theorembox}

\begin{intuitionbox}{La Delta como Filtro / Interruptor de Sustitución}
Cuando la delta de Kronecker se contrae con las componentes de un vector $A_j$, actúa como un **filtro selector**:
\[
\delta_{ij} A_j = \delta_{i1} A_1 + \delta_{i2} A_2 + \delta_{i3} A_3.
\]
De los tres términos de la suma, dos se apagan a cero inmediatamente. El único término sobreviviente es aquel donde el índice mudo $j$ coincide con el índice libre $i$. Por lo tanto:
\[
\mathbf{\delta_{ij} A_j = A_i}.
\]
\textbf{Regla operativa:} Contraer $\delta_{ij}$ con cualquier tensor consiste simplemente en borrar la delta y reemplazar el índice repetido $j$ por el índice $i$.
\end{intuitionbox}

\begin{examplebox}{Cálculo de la Traza de la Delta}
¿Cuánto vale la contracción consigo misma $\delta_{ii}$?
\[
\delta_{ii} = \delta_{11} + \delta_{22} + \delta_{33} = 1 + 1 + 1 = 3.
\]
En un espacio de dimensión $n$, la contracción $\delta_{ii} = n$ mide la dimensión del espacio. Asimismo:
\[
\delta_{ij}\delta_{jk} = \delta_{ik}, \qquad \delta_{ij}\delta_{ij} = \delta_{ii} = 3.
\]
\end{examplebox}

% =====================================================================
\manualsection[4. Levi-Civita $\epsilon_{ijk}$]{4. El Símbolo de Levi-Civita $\epsilon_{ijk}$: Geometría de Permutaciones}

El producto cruz entre dos vectores es una operación puramente antisimétrica: $\vec{A} \times \vec{B} = -\vec{B} \times \vec{A}$, lo que implica que $\vec{A} \times \vec{A} = \vec{0}$. Para codificar esta alternancia introducimos el \term{símbolo de Levi-Civita} en $\R^3$.

\begin{theorembox}{Definición del Símbolo de Levi-Civita}
El símbolo $\epsilon_{ijk}$ (con $i, j, k \in \{1, 2, 3\}$) es un objeto con $3^3 = 27$ componentes definido por:
\[
\epsilon_{ijk} = \begin{cases} 
+1 & \text{si } (i, j, k) \text{ es una permutación par (cíclica) de } (1, 2, 3), \\
-1 & \text{si } (i, j, k) \text{ es una permutación impar (anticíclica) de } (1, 2, 3), \\
0 & \text{si dos o más índices son iguales.}
\end{cases}
\]
\end{theorembox}

\begin{center}
\begin{tikzpicture}[scale=1.1]
    % Clock circle
    \draw[line width=1.5pt, color=orangeacc] (0,0) circle (1.6cm);
    
    % Nodes 1, 2, 3
    \node[circle, fill=blackmain, text=white, font=\bfseries\large, inner sep=4pt] (N1) at (90:1.6cm) {1};
    \node[circle, fill=blackmain, text=white, font=\bfseries\large, inner sep=4pt] (N2) at (330:1.6cm) {2};
    \node[circle, fill=blackmain, text=white, font=\bfseries\large, inner sep=4pt] (N3) at (210:1.6cm) {3};
    
    % Direction arrows (clockwise = +1)
    \draw[-{Stealth[length=3mm]}, line width=1.8pt, color=orangeacc] (70:1.9cm) arc (70:10:1.9cm) node[midway, right=2pt, font=\bfseries\small, color=orangeacc] {$+1$ (horario)};
    \draw[-{Stealth[length=3mm]}, line width=1.8pt, color=orangeacc] (310:1.9cm) arc (310:250:1.9cm);
    \draw[-{Stealth[length=3mm]}, line width=1.8pt, color=orangeacc] (190:1.9cm) arc (190:130:1.9cm);

    % Legend
    \node[align=center, font=\small] at (0, 0) {\textbf{Permutación}\\\textbf{Cíclica}};
\end{tikzpicture}
\end{center}

De las 27 componentes posibles:
\begin{itemize}
    \item **Solo 3 componentes valen $+1$** (sentido horario):
    \[
    \epsilon_{123} = \epsilon_{231} = \epsilon_{312} = +1.
    \]
    \item **Solo 3 componentes valen $-1$** (sentido antihorario, intercambio de un par):
    \[
    \epsilon_{132} = \epsilon_{213} = \epsilon_{321} = -1.
    \]
    \item **Las 21 componentes restantes valen $0$** (por ejemplo, $\epsilon_{112} = \epsilon_{323} = 0$).
\end{itemize}

\begin{theorembox}{El Producto Cruz en Notación de Índices}
La $i$-ésima componente del producto cruz entre $\vec{A}$ y $\vec{B}$ es exactamente:
\[
\mathbf{(\vec{A} \times \vec{B})_i = \epsilon_{ijk} A_j B_k}.
\]
\end{theorembox}

\begin{examplebox}{Comprobación para la Componente $x$ ($i = 1$)}
Evaluemos la primera componente según la fórmula:
\[
(\vec{A} \times \vec{B})_1 = \epsilon_{1jk} A_j B_k = \epsilon_{123} A_2 B_3 + \epsilon_{132} A_3 B_2 = (+1)A_2 B_3 + (-1)A_3 B_2 = A_2 B_3 - A_3 B_2.
\]
¡Recuperamos exactamente $A_y B_z - A_z B_y$ sin necesidad de dibujar determinantes $3\times 3$!
\end{examplebox}

% =====================================================================
\manualsection[5. La Identidad Maestra]{5. La Identidad Maestra $\epsilon-\delta$: El Puente Fundamental}

El poder absoluto del álgebra de índices reside en contraer dos símbolos de Levi-Civita que comparten un índice común.

\begin{theorembox}{Identidad Maestra de Contracción $\epsilon-\delta$}
Al sumar sobre el primer índice compartido $i$, se cumple la identidad universal:
\[
\mathbf{\epsilon_{ijk}\epsilon_{imn} = \delta_{jm}\delta_{kn} - \delta_{jn}\delta_{km}}.
\]
\end{theorembox}

\begin{intuitionbox}{Mnemotecnia Visual para Recordar la Identidad}
Observa con atención el orden de los índices en el segundo miembro:
\[
\epsilon_{i\mathbf{j}\mathbf{k}}\epsilon_{i\mathbf{m}\mathbf{n}} = 
\delta_{\mathbf{j}\mathbf{m}}\delta_{\mathbf{k}\mathbf{n}} - \delta_{\mathbf{j}\mathbf{n}}\delta_{\mathbf{k}\mathbf{m}}.
\]
\begin{itemize}
    \item El término positivo empareja **los primeros con los primeros** ($j$ con $m$) y **los segundos con los segundos** ($k$ con $n$).
    \item El término negativo invierte el emparejamiento con un signo menos cruzado: **los extremos** ($j$ con $n$) y **los internos** ($k$ con $m$).
\end{itemize}
Dado que el símbolo $\epsilon_{ijk}$ es invariante bajo permutaciones cíclicas ($\epsilon_{ijk} = \epsilon_{jki} = \epsilon_{kij}$), si el índice compartido se encuentra en otra posición, basta rotarlo cíclicamente para aplicar la fórmula.
\end{intuitionbox}

% =====================================================================
\manualsection[6. Laboratorio de Demolición]{6. Laboratorio de Demolición: Identidades Vectoriales en 3 Líneas}

Pongamos a prueba el arsenal demostrando las identidades vectoriales más famosas de la física matemática.

\manualsubsection{Demostración 1: La Identidad BAC--CAB}
Queremos probar que $\vec{A} \times (\vec{B} \times \vec{C}) = \vec{B}(\vec{A}\cdot\vec{C}) - \vec{C}(\vec{A}\cdot\vec{B})$.

\begin{examplebox}{Demostración en 3 Renglones}
Analizamos la $i$-ésima componente del miembro izquierdo:
\begin{align*}
[\vec{A} \times (\vec{B} \times \vec{C})]_i &= \epsilon_{ijk} A_j (\vec{B} \times \vec{C})_k \\
&= \epsilon_{ijk} A_j \big(\epsilon_{kmn} B_m C_n\big) = \epsilon_{kij}\epsilon_{kmn} A_j B_m C_n.
\end{align*}
Aplicamos la identidad maestra $\epsilon-\delta$ sobre el índice compartido $k$:
\[
\epsilon_{kij}\epsilon_{kmn} = \delta_{im}\delta_{jn} - \delta_{in}\delta_{jm}.
\]
Sustituyendo y contrayendo con las deltas:
\begin{align*}
[\vec{A} \times (\vec{B} \times \vec{C})]_i &= (\delta_{im}\delta_{jn} - \delta_{in}\delta_{jm}) A_j B_m C_n \\
&= \delta_{im} B_m (\delta_{jn} A_j C_n) - \delta_{in} C_n (\delta_{jm} A_j B_m) \\
&= B_i (A_j C_j) - C_i (A_j B_j).
\end{align*}
Reconociendo los productos escalares $A_j C_j = \vec{A}\cdot\vec{C}$ y $A_j B_j = \vec{A}\cdot\vec{B}$:
\[
[\vec{A} \times (\vec{B} \times \vec{C})]_i = B_i (\vec{A}\cdot\vec{C}) - C_i (\vec{A}\cdot\vec{B}) = [\vec{B}(\vec{A}\cdot\vec{C}) - \vec{C}(\vec{A}\cdot\vec{B})]_i \quad \blacksquare
\]
\end{examplebox}

\manualsubsection{Demostración 2: La Identidad de Lagrange}
Demostremos la relación entre las normas del producto cruz y el producto punto:
\[
|\vec{A} \times \vec{B}|^2 = |\vec{A}|^2|\vec{B}|^2 - (\vec{A}\cdot\vec{B})^2.
\]

\begin{examplebox}{Demostración en 2 Renglones}
Por definición de norma euclidiana:
\[
|\vec{A} \times \vec{B}|^2 = (\vec{A} \times \vec{B})_i (\vec{A} \times \vec{B})_i = (\epsilon_{ijk} A_j B_k)(\epsilon_{imn} A_m B_n) = \epsilon_{ijk}\epsilon_{imn} A_j B_k A_m B_n.
\]
Usando directamente la identidad $\epsilon_{ijk}\epsilon_{imn} = \delta_{jm}\delta_{kn} - \delta_{jn}\delta_{km}$:
\begin{align*}
|\vec{A} \times \vec{B}|^2 &= (\delta_{jm}\delta_{kn} - \delta_{jn}\delta_{km}) A_j B_k A_m B_n \\
&= (A_j A_j)(B_k B_k) - (A_j B_j)(B_k A_k) \\
&= |\vec{A}|^2 |\vec{B}|^2 - (\vec{A}\cdot\vec{B})^2 \quad \blacksquare
\end{align*}
\end{examplebox}

% =====================================================================
\manualsection[7. Aplicaciones en Física]{7. Aplicaciones en Física: Momento Angular y Fuerza de Lorentz}

\manualsubsection{1. El Teorema de Simetría vs. Antisimetría}
Una de las propiedades más poderosas en física teórica es la contracción entre un tensor simétrico y uno antisimétrico.

\begin{theorembox}{Anulación de Contracción Simétrico--Antisimétrico}
Si $S_{ij} = S_{ji}$ es simétrico y $A_{ij} = -A_{ji}$ es antisimétrico, entonces:
\[
\mathbf{S_{ij} A_{ij} = 0}.
\]
\end{theorembox}
*Demostración:* Como los índices $i$ y $j$ son mudos, podemos renombrarlos $i \leftrightarrow j$:
\[
S_{ij} A_{ij} = S_{ji} A_{ji} = (S_{ij})(-A_{ij}) = -S_{ij} A_{ij} \implies 2 S_{ij} A_{ij} = 0 \implies S_{ij} A_{ij} = 0.
\]

\manualsubsection{2. ¿Por qué el campo magnético no realiza trabajo?}
La fuerza de Lorentz magnética sobre una carga $q$ con velocidad $\vec{v}$ es:
\[
F_i = q\,\epsilon_{ijk} v_j B_k.
\]
La potencia instantánea transmitida a la partícula es el producto escalar de la fuerza y la velocidad:
\[
P = \vec{F}\cdot\vec{v} = F_i v_i = q\,\epsilon_{ijk} v_i v_j B_k.
\]
Notemos que el producto $v_i v_j$ es perfectamente **simétrico** bajo el intercambio de $i$ y $j$ ($v_i v_j = v_j v_i$), mientras que $\epsilon_{ijk}$ es **antisimétrico** ($\epsilon_{jik} = -\epsilon_{ijk}$). Por el teorema anterior:
\[
\mathbf{P = 0}.
\]
El campo magnético desvía la trayectoria sin alterar la energía cinética de la partícula. ¡Demostrado rigurosamente en un solo renglón!

\manualsubsection{3. Momento Angular y Torque Central}
El momento angular de una partícula puntual es $\vec{L} = \vec{r}\times\vec{p}$, y el torque de una fuerza es $\vec{\tau} = \vec{r}\times\vec{F}$:
\[
L_i = \epsilon_{ijk} x_j p_k, \qquad \tau_i = \epsilon_{ijk} x_j F_k.
\]
Si la fuerza es central, $\vec{F} = f(r)\hat{r} = \frac{f(r)}{r}\vec{r}$, sus componentes son $F_k = g(r) x_k$. Entonces:
\[
\tau_i = \epsilon_{ijk} x_j \big(g(r) x_k\big) = g(r)\,\epsilon_{ijk} x_j x_k = 0 \quad (\text{pues } x_j x_k \text{ es simétrico}).
\]
Por lo tanto, $\frac{d\vec{L}}{dt} = \vec{\tau} = \vec{0}$, garantizando la conservación del momento angular en campos gravitacionales o electrostáticos.

\end{document}
